\documentclass[11pt]{article}
\usepackage[margin=1in]{geometry}
\usepackage{amsmath,amssymb,amsthm,booktabs}
\newtheorem{theorem}{Theorem}
\title{Permanent and Determinant Congruences in $3\times3$ Digit Grids}
\author{Michael Warren Song\\ \small computed with Egor; reproduction: \texttt{reproduce.py} (plain Python 3)}
\date{November 2026}
\begin{document}
\maketitle

\section*{Status}
\textbf{Proved:} Theorem~1. \textbf{Executed:} every number below, by \texttt{reproduce.py}, which prints PASS/FAIL per claim.
The criterion has not been checked against the literature, so no claim is made that it is new.

\section{The congruence}
For a $3\times3$ matrix, $\operatorname{per}(A)-\det(A)=2\,O(A)$, where $O(A)$ is the sum of the three products along odd
permutations. Since 2 is invertible modulo 9, $\operatorname{per}(A)\equiv\det(A)\pmod 9$ exactly when $O(A)\equiv0\pmod9$.

\begin{theorem}
Let $A_{ij}=u_i+3v_j$, where $u_1,u_2,u_3$ lie in the three distinct classes modulo 3. Then
\[\operatorname{per}(A)\equiv\det(A)\pmod 9\iff v_1+v_2+v_3\equiv0\pmod3.\]
\end{theorem}
\begin{proof}
Let $r$ be the row with $3\mid u_r$, and $a,b$ the other two rows. Each of the three odd permutations (the transpositions)
sends $r$ to a different column, and each factor $A_{rj}=3(u_r/3+v_j)$ is divisible by 3. Every entry of rows $a$ and $b$ is
$\equiv u_a$ or $u_b\pmod 3$, whatever its column. Hence
\[O(A)\equiv 3\,u_au_b\sum_{j=1}^{3}\Bigl(\tfrac{u_r}{3}+v_j\Bigr)=3\,u_au_b\,(u_r+v_1+v_2+v_3)\equiv 3\,u_au_b\,(v_1+v_2+v_3)\pmod 9,\]
using $3\mid u_r$. As $u_au_b\not\equiv0\pmod3$, $O(A)\equiv0\pmod9$ if and only if $v_1+v_2+v_3\equiv0\pmod3$.
\end{proof}
The grid $\begin{smallmatrix}1&2&3\\4&5&6\\7&8&9\end{smallmatrix}$ is the transpose of the case $u=(1,2,3)$, $v=(0,1,2)$, and
$0+1+2=3$, so its permanent and determinant agree modulo 9. A direct check over $u_i,v_j\in\{0,\dots,8\}$ (118\,098 matrices)
finds no exception.

\section{The two swaps}
\begin{center}
\begin{tabular}{@{}lrrr@{}}
\toprule
swap in $1\,2\,3/4\,5\,6/7\,8\,9$ & permanent & determinant & odd-permutation sum\\
\midrule
none & 450 & 0 & 225\\
$2\leftrightarrow5$ & 594 & 54 & 270\\
$7\leftrightarrow8$ & 459 & $-9$ & 234\\
both & 612 & 72 & 270\\
\bottomrule
\end{tabular}
\end{center}
All four agree modulo 9. The determinant's interaction term is $72-54-(-9)+0=27$, and the two single swaps give
determinants of opposite sign, so the parity of a swap does not fix the sign of the result.

\section{Row flips}
Reversing rows of $1\,2\,3/4\,5\,6/7\,8\,9$ (the flips of \texttt{all\_flips\_123\_456\_789\_2026\_10\_19.py}) keeps the
determinant 0, since every row stays evenly spaced and so in the span of $(1,1,1)$ and $(1,2,3)$. The permanent takes the
values 450, 478, 490, 502. It agrees with the determinant modulo 9 only with no flip or all three flips; one or two flips
leave remainders 1, 4 or 7. A flipped row leaves the family of Theorem~1.

\section*{Scope}
These are statements about fixed matrices. They say nothing about the circuit size needed to compute the permanent as the
size grows (Valiant's conjecture).
\end{document}
